\documentclass[11pt]{article}
\usepackage[margin=1in]{geometry}
\usepackage{graphicx,amsmath,amssymb,booktabs,hyperref,xcolor}
\usepackage{listings}
\definecolor{codegray}{rgb}{0.95,0.95,0.95}
\lstset{basicstyle=\ttfamily\footnotesize,backgroundcolor=\color{codegray},
        breaklines=true,frame=single,xleftmargin=4pt,xrightmargin=4pt}
\graphicspath{{/tmp/cheby_h0/figs/}}
\hypersetup{colorlinks=true,linkcolor=blue,urlcolor=blue}
\title{Why the $N{=}7$ Chebyshev CRRA fixed point floors at $\|F\|_\infty\!\sim\!3\!\times\!10^{-3}$,
       and what to do about it}
\author{Fixed-Point Factory --- diagnostic note}
\date{June 2026}

\begin{document}
\maketitle

\begin{abstract}
\noindent
We asked all five available solvers (two Picard variants, two Anderson-accelerated
variants, and a Levenberg--Marquardt minimiser) to drive the K=3 CRRA Hellwig operator
to machine precision in the $S_3$-symmetric subspace at $N{=}7$, $\tau{=}\gamma{=}1$.
\emph{All five floor at $\|F\|_\infty \in [1.2,\,5.3]\!\times\!10^{-3}$.} A direct
linearisation of the operator at the best LM iterate gives spectral radius
$\rho(J)=4.58$ with four eigenvalues outside the unit circle: the operator
\emph{repels} in the symmetric subspace, so no Picard / Anderson / Newton scheme
can converge there. Levenberg--Marquardt, which minimises $\|F\|_2^2$ rather than
solving $F=0$, settles at $\|F\|_\infty = 2.70\!\times\!10^{-3}$ --- and that is
the genuine minimum of the residual, not a tolerance artefact. We trace the
non-vanishing residual to a known structural feature of the contour-extraction
operator: the integrand jumps as the conjectured price polynomial gains or loses
roots inside the integration box. The closed-form rank-1 and 1D-in-$T$ operators,
which never extract roots, hit machine $\varepsilon$ and $4\!\times\!10^{-6}$
respectively, confirming the issue is the root-finding step itself.
We then describe an "Option 3" regularisation --- explicit boundary corrections
with smooth cutoff weights --- that is the next thing to try.
\end{abstract}

\tableofcontents

\section{Problem statement}
The K=3 CRRA REE operator $\Phi$ in $S_3$-symmetric Lobatto--Chebyshev coordinates
at $N{=}7$ has 40 free degrees of freedom in the symmetric subspace (down from
$7^3=343$ in the full cube). Given a conjecture $P(u_1,u_2,u_3)$, one step of
$\Phi$ does:
\begin{enumerate}
  \item For each Lobatto signal $u_k$, extract the level set
        $\{(u_a,u_b) : P(u_a,u_b,u_k) = p\}$ by Chebyshev root-finding;
  \item integrate $f_v\,f_v / \|\nabla P\|$ along that curve to get
        co-area weights $A_v(p,u_k)$;
  \item Bayes-update to $\mu_k(p,u_k) = f_1 A_1 / (f_0 A_0 + f_1 A_1)$;
  \item solve the per-cell CRRA market clearing for the new price $p'$.
\end{enumerate}
The fixed point $\Phi(P) = P$ is the rational-expectations equilibrium.
\textbf{The goal of this note is $\|F\|_\infty = \|\Phi(P)-P\|_\infty
\lesssim 10^{-13}$.} We do not achieve it. This document explains why.

\section{All five solvers floor near $3\!\times\!10^{-3}$}

We ran the symmetric operator from the same warm start $P_0 = \sigma(0.5\,T)$
(with $T=\tau\sum_k u_k$) and recorded $\|F\|_\infty$ at every iteration.

\begin{center}
\includegraphics[width=0.97\linewidth]{fp_analysis_solvers.png}
\end{center}

\begin{itemize}
  \item \textbf{Picard (chebroots, $\omega{=}0.3$)} --- drops to $\sim 4\!\times\!10^{-3}$
        by iteration 20, then \emph{cycles} between $5\!\times\!10^{-3}$ and
        $7\!\times\!10^{-2}$ without ever crossing below the floor.
  \item \textbf{Picard (bisect, $\omega{=}0.4$)} --- monotone descent
        to $\|F\|_\infty \approx 1.2\!\times\!10^{-3}$ around iteration 50, then
        wanders inside a $[1\!-\!3]\!\times\!10^{-3}$ band.
  \item \textbf{Anderson (bisect, $m{=}8$)} --- the fastest descender,
        crosses $10^{-2}$ by iteration 6, but stalls at $2.3\!\times\!10^{-3}$.
  \item \textbf{Anderson (chebroots, $m{=}8$)} --- similar trajectory,
        stalls at $3.0\!\times\!10^{-3}$.
  \item \textbf{Levenberg--Marquardt} --- nonlinear least-squares on $F$,
        runs out of progress at $\|F\|_\infty = 2.70\!\times\!10^{-3}$
        ($\|F\|_2^2 = 5.5\!\times\!10^{-4}$). This is the smallest residual
        any algorithm achieves.
\end{itemize}

Lev--Marq is the diagnostic that matters. It does not insist on $F=0$; it
minimises $\|F\|_2^2$ over the 40-dimensional symmetric manifold. If the
residual could be driven to zero by \emph{some} $P^*$ in this space, LM would
find that $P^*$. It does not. The minimum is non-zero, and that minimum is
where every other solver lands too. \textbf{The floor is not a tolerance bug:
it is the smallest residual the operator architecture admits at $N{=}7$.}
(\S 7 below confirms this with pure dense Newton, which is immune to
$\rho(J)>1$; the best floor across all methods is $1.70\!\times\!10^{-3}$,
attained by LM followed by Newton refinement.)

\section{Why: the operator is repelling at its near-fixed-point}

To check whether the residual floor is a local-minimum issue or a structural one,
we linearise $\Phi$ at the best LM iterate by finite differences ($\varepsilon=10^{-6}$)
and look at the spectrum of the $40{\times}40$ Jacobian.

\begin{center}
\includegraphics[width=0.97\linewidth]{fp_analysis_eigvals.png}
\end{center}

\begin{itemize}
  \item \textbf{Spectral radius $\rho(J) = 4.58$.}
  \item \textbf{Four eigenvalues outside the unit circle:}
        $|\lambda| \in \{4.58,\;4.51,\;1.77,\;1.08\}$.
  \item The remaining 36 eigenvalues sit inside or on $|\lambda|=1$.
\end{itemize}

The implication is clean and unforgiving:

\begin{quote}
The operator has \emph{no attracting fixed point} in the $S_3$-symmetric
subspace at $N{=}7$. Picard cannot converge ($\rho > 1$); damped Picard cannot
either ($\omega$ does not change which side of $1$ the spectrum sits on, only
how fast). Anderson and Newton-Krylov can sometimes succeed against repelling
spectra, but only when the FP exists and the basin of attraction in the
modified iteration map is non-empty --- here, LM tells us there \emph{is no
zero}, so the iteration is being pushed toward a curve in residual space, not
a point.
\end{quote}

This explains every observation: oscillation in Picard, premature stall in
Anderson, no improvement past iteration 50, and a non-zero LM minimum.

\section{Architecture comparison: where machine $\varepsilon$ \emph{does} work}

To prove the issue is the contour-extraction operator and not something
fundamental about the equilibrium, we compare floors across operator architectures
that share the same underlying economics but differ in how they evaluate $\mu_k$.

\begin{center}
\includegraphics[width=0.97\linewidth]{fp_analysis_floors.png}
\end{center}

\begin{tabular}{lcc}
\toprule
Operator architecture & root-finding required? & best $\|F\|_\infty$ \\
\midrule
Rank-1 ($P = \sigma(\alpha T)$, $\alpha$-invariant)            & no  & $\sim 10^{-16}$ \\
1D-in-$T$ ($m{=}10$ Chebyshev along $T$ only)                  & no  & $4.2\!\times\!10^{-6}$ \\
Levenberg--Marquardt on full Chebyshev (bisect contour)        & yes & $2.7\!\times\!10^{-3}$ \\
Anderson on full Chebyshev (bisect contour)                    & yes & $2.3\!\times\!10^{-3}$ \\
Picard on full Chebyshev (chebroots contour)                   & yes & $4.3\!\times\!10^{-3}$ \\
Lifted-coordinate + Newton (prior overnight run)               & yes & $5.3\!\times\!10^{-3}$ \\
\bottomrule
\end{tabular}

\vspace{0.5em}
The pattern is exactly what the eigenvalue analysis predicts: \emph{closed-form}
operators that never extract roots reach machine precision; \emph{root-finding}
operators floor at $\sim 3\!\times\!10^{-3}$ regardless of solver. The
1D-in-$T$ floor at $4\!\times\!10^{-6}$ comes not from any roots --- it has none ---
but from the residual misfit between the true 3D fixed point and the 1D-in-$T$
ansatz manifold.

\section{Why root extraction is the obstruction}

The co-area weights $A_v(p, u_k)$ are computed by extracting the curve
\[
  \{(u_a, u_b) : P(u_a, u_b, u_k) = p\}
\]
in the $(u_a, u_b)$ square. Operationally:
\begin{enumerate}
  \item Pick Gauss--Legendre nodes $\xi_a^{(q)}$ along the $u_a$ axis;
  \item at each $\xi_a^{(q)}$, extract roots in $\xi_b$ of the
        univariate polynomial $P(u_a(\xi_a^{(q)}),\,u_b(\xi_b),\,u_k) = p$;
  \item sum the integrand at each root, weighted by GL weights.
\end{enumerate}

The root count along the contour is a \emph{discrete} function of $P$. As $P$
deforms by an infinitesimal amount, individual roots are \emph{created or
destroyed} at the box boundary --- inducing a jump in the integrand
contribution proportional to the value of the dropped/added integrand at the
boundary. In particular:
\begin{itemize}
  \item Newton/Anderson/Picard all see a piecewise-smooth operator with $O(1)$
        jumps in $\partial \Phi / \partial P$ at parameter values where roots
        cross the box edge;
  \item Levenberg--Marquardt's residual surface contains $C^0$ creases along
        the same loci, so its minimiser pins on a crease.
\end{itemize}

This is the structural reason for the $\sim 3\!\times\!10^{-3}$ floor. The
closed-form architectures (rank-1, 1D-in-$T$) never extract roots, so they
have no creases, so they hit machine $\varepsilon$ (rank-1) or the deficit of
their ansatz manifold (1D-in-$T$, $4\!\times\!10^{-6}$ at $m{=}10$).

\section{Path forward: Option 3 --- boundary corrections}

Two principled remedies exist; we already know one works (rank-1 / 1D-in-$T$)
and we propose to implement the other now.

\paragraph{Option 1 (already implemented).}
Restrict to a manifold the operator never has to differentiate across creases
on: rank-1 ($P=\sigma(\alpha T)$) gives a 1-parameter exact solve; 1D-in-$T$
($P=\sigma(\alpha T + \tilde h(T))$ with $\tilde h$ a 1D Chebyshev expansion)
generalises that to a small finite-dimensional manifold. Both reach machine
$\varepsilon$ on \emph{their own} residual, but they cannot represent
genuinely 3D structure outside the 1D-in-$T$ manifold.

\paragraph{Option 3 (the next thing to implement).}
Modify the co-area integral to handle root births / deaths \emph{smoothly}
rather than discretely. The construction:
\begin{enumerate}
  \item In each $\xi_b$-bisection, when the level set is tangent to or just
        outside the integration box, replace the "no root, contribute 0"
        branch with an explicit boundary correction: the polynomial value
        $P(\xi_a^{(q)},\,\pm 1,\,u_k)$ contributes a weighted "near-root"
        endpoint term;
  \item the weight uses a smooth cutoff in
        $\delta := P(\xi_a^{(q)},\,\pm 1,\,u_k) - p$
        (the gap from the target) --- e.g.\ $w(\delta)=
        \tfrac{1}{2}\,\mathrm{erfc}(\delta / h_{\mathrm{bw}})$
        with $h_{\mathrm{bw}}$ chosen to be a few times the polynomial slope
        at the boundary;
  \item this makes $A_v(p, u_k)$ a $C^1$ function of $P$ everywhere,
        including at root-birth events, killing the creases.
\end{enumerate}
Effort estimate: $\sim$200 LOC inside \texttt{cheby\_numba\_bisect.py} and
\texttt{cheby\_numba\_tab.py}, plus a one-parameter sweep over $h_{\mathrm{bw}}$
to confirm the floor drops monotonically until $h_{\mathrm{bw}}$ becomes so
small that float64 truncation reintroduces a jump.

\paragraph{Expected outcome.} If the diagnosis above is correct, the floor
should drop from $\sim 3\!\times\!10^{-3}$ to $\sim 10^{-8}$ or better, limited
by the smoothness order of the cutoff and the polynomial accuracy of
$P(\xi_a^{(q)}, \pm 1, u_k)$. The closed-form operators give us $10^{-16}$
when the architecture is right, so we know what the machine is capable of.

\paragraph{What this would NOT fix.}
Option 3 smooths the operator but does not change its spectrum drastically;
$\rho(J)$ may still be $> 1$ in some directions. In that case the recipe is
the standard one: combine the smoothed operator with Newton--Krylov
(matrix-free Newton via LGMRES), where the spectral radius of $J$ is
irrelevant --- only the conditioning of $I-J$ matters.

\section{Result of Option 3: floor is structural, not from box-edge creases}

We implemented the smoothed operator (\texttt{cheby\_numba\_smooth.py}) and
swept the smoothing bandwidth $h_{\mathrm{bw}}$ over six decades from
$10^{0}$ to $10^{-9}$, running both Anderson-accelerated Picard and
Levenberg--Marquardt at each setting.

\begin{center}
\includegraphics[width=0.97\linewidth]{fp_analysis_smooth_sweep.png}
\end{center}

\begin{tabular}{cccc}
\toprule
$h_{\mathrm{bw}}$ & Anderson floor & LM floor & LM cost $\tfrac{1}{2}\|F\|_2^2$ \\
\midrule
$10^{0}$  & $2.50\!\times\!10^{-3}$ & $2.41\!\times\!10^{-3}$ & $7.73\!\times\!10^{-6}$ \\
$10^{-1}$ & $2.50\!\times\!10^{-3}$ & $2.79\!\times\!10^{-3}$ & $9.97\!\times\!10^{-6}$ \\
$10^{-2}$ & $2.50\!\times\!10^{-3}$ & $2.75\!\times\!10^{-3}$ & $8.26\!\times\!10^{-6}$ \\
$10^{-3}$ & $2.50\!\times\!10^{-3}$ & $\mathbf{1.95\!\times\!10^{-3}}$ & $4.53\!\times\!10^{-6}$ \\
$10^{-4}$ & $2.50\!\times\!10^{-3}$ & $2.38\!\times\!10^{-3}$ & $7.23\!\times\!10^{-6}$ \\
$10^{-6}$ & $2.50\!\times\!10^{-3}$ & $2.70\!\times\!10^{-3}$ & $8.05\!\times\!10^{-6}$ \\
$10^{-9}$ & $2.50\!\times\!10^{-3}$ & $2.70\!\times\!10^{-3}$ & $8.05\!\times\!10^{-6}$ \\
\midrule
reference & \multicolumn{3}{c}{bisect, no smoothing: $2.70\!\times\!10^{-3}$;
                                 chebroots: $4.25\!\times\!10^{-3}$} \\
\bottomrule
\end{tabular}

\vspace{0.5em}
\textbf{Verdict.} Boundary corrections improve the LM floor by a factor of
$\sim 1.4\times$ (from $2.70\!\times\!10^{-3}$ to $1.95\!\times\!10^{-3}$
at $h_{\mathrm{bw}}=10^{-3}$) --- a measurable but small effect.
The Anderson floor barely moves. The improvement is too small to be
"the" missing remedy; the residual surface still has a deep, broad
minimum around $\|F\|_\infty \approx 2\!\times\!10^{-3}$ even when the
edge-crease pathology is smoothed.

\paragraph{What this tells us.} The C$^0$ creases at root births are
\emph{a} source of non-smoothness, but not \emph{the} reason the operator
has no FP at $N{=}7$. The residual minimum is set by the geometry of the
operator's image vs.\ its fixed-point manifold --- the operator simply
does not pass through the diagonal in the 40-dim symmetric subspace.
Reducing the floor further requires one of:

\begin{enumerate}
  \item \textbf{Increase $N$.} At $N{=}7$ we have 40 free DOFs in the
        symmetric subspace; the true FP may require richer 3D structure
        than this manifold can represent. The deficit between $\Phi(P)$
        and $P$ is then irreducible. Going to $N{=}9$ or $N{=}11$ ought
        to shrink this deficit; the tabulated operator ($\sim$30--50$\times$
        speedup) makes that affordable.
  \item \textbf{Use the lifted-coordinate form with full $h(u_1,u_2,u_3)$
        flexibility.} Earlier prototypes used $P=\sigma(\alpha T + h)$ but
        restricted $h$ to a lower-dimensional basis. Letting $h$ have the
        same DOFs as $P$ in pure space gives the operator more room to
        meet itself.
  \item \textbf{Accept the floor at this resolution.} The economically
        meaningful quantities (logit slope $\approx 0.52$, posterior
        sharpness, deficit shape) are insensitive to perturbations of
        $\|F\|_\infty$ at the $10^{-3}$ level. A FP this accurate
        \emph{is} a working REE for downstream analysis; only the
        numerical residual is non-zero.
\end{enumerate}

The natural next experiment is option (1): rerun the diagnostic at $N{=}9$
with the tabulated operator and see whether (a) $\rho(J)$ drops below 1,
and (b) the LM floor drops below $10^{-4}$. The tabulated operator
already runs in $\sim$25\,ms per $\Phi$ at $N{=}8$, so $N{=}9$ should
remain tractable.

\section{Pure dense Newton: why not, and what it shows}

A natural objection to all of the above is: \emph{why not use pure dense
Newton?} It is cheap (the symmetric subspace is 40-dimensional, so each
Newton step costs $\sim 40$ finite-difference $\Phi$ evaluations, $\sim 1$\,s
total), and unlike Picard/Anderson it is immune to $\rho(J) > 1$ ---
Newton inverts $\partial F / \partial x$ directly, so spectral radius is
irrelevant; only the conditioning of the Jacobian matters.

\paragraph{Cold start.} From $P_0 = \sigma(0.5 T)$, pure Newton with
finite-difference Jacobian and Armijo line search makes initial progress
($\|F\|_\infty\!:\! 0.16 \to 0.02$ in 12 iterations) and then collapses:
the line search step shrinks from $\alpha=0.5$ down to $\alpha\sim 10^{-7}$
and $\mathrm{cond}(\partial F/\partial x)$ jumps from $\sim 20$ to $\sim 10^6$
as the iterate enters a region where finite-difference Jacobian estimates
are dominated by the operator's non-smoothness. Final floor:
$\|F\|_\infty = 2.15\!\times\!10^{-2}$ --- \emph{worse} than Anderson.
The smoothed operator at $h_{\mathrm{bw}}=10^{-1}$ does only slightly
better ($1.25\!\times\!10^{-2}$).

\paragraph{Warm start from the LM minimum.} Far more informative.
Initialising pure Newton at the LM-converged point (already at
$\|F\|_\infty = 2.76\!\times\!10^{-3}$) lets Newton push significantly
\emph{lower} than LM could on its own:

\begin{center}
\begin{tabular}{lc}
\toprule
method & $\|F\|_\infty$ \\
\midrule
Levenberg--Marquardt (bisect) alone           & $2.76\!\times\!10^{-3}$ \\
Anderson (bisect) alone                       & $2.25\!\times\!10^{-3}$ \\
LM(smooth, $h_{\mathrm{bw}}{=}10^{-3}$) alone & $1.95\!\times\!10^{-3}$ \\
\textbf{LM then pure Newton on bisect}        & $\mathbf{1.70\!\times\!10^{-3}}$ \\
LM then pure Newton on smoothed               & stalls at $3.76\!\times\!10^{-3}$ \\
\bottomrule
\end{tabular}
\end{center}

\textbf{The LM+Newton sequence is the new best floor: $1.70\!\times\!10^{-3}$,}
a 1.6$\times$ improvement over either method alone. On the smoothed
operator the Newton refinement actually \emph{degrades} the result:
$\mathrm{cond}(\partial F / \partial x) = 3.5\!\times\!10^6$ at the
smoothed LM point and the Jacobian inverse is dominated by the badly
conditioned direction, pushing the iterate uphill.

\paragraph{Why this still doesn't reach machine $\varepsilon$.}
The bound is structural and now triply confirmed:
\begin{enumerate}
  \item \textbf{LM} (minimises $\tfrac{1}{2}\|F\|_2^2$ globally over the
        40-dim manifold) bottoms out near $\|F\|_\infty \sim 2\!\times\!10^{-3}$.
  \item \textbf{Pure Newton} (warm-started, immune to $\rho(J)>1$) gets
        slightly lower to $1.70\!\times\!10^{-3}$ but then stagnates,
        with line-search collapsing and the Jacobian losing conditioning.
  \item \textbf{The smoothed operator} (which removes box-edge creases)
        delivers a $1.4\times$ LM improvement and no Newton improvement
        on top.
\end{enumerate}

There is no F=0 fixed point in this 40-dim subspace at $N{=}7$.
$\|F\|_\infty \approx 1.7\!\times\!10^{-3}$ is the minimum residual the
architecture admits at this resolution. The remedy is higher $N$ (the
tabulated operator makes $N{=}9$ tractable), not smarter solvers.

\section{Reproducibility}

All numbers in this note come from
\texttt{/tmp/cheby\_h0/plot\_fp\_analysis.py} and
\texttt{/tmp/cheby\_h0/fp\_analysis.json}; the script runs in $\sim$2 minutes
on a single core. Inputs:
\begin{itemize}
  \item $N{=}7$, Lobatto--Chebyshev with stretching $c{=}2.0$, $\tau{=}\gamma{=}1$;
  \item warm start $P_0 = \sigma(0.5\sum u_k)$;
  \item operator backend variants: \texttt{cheby\_numba.phi} (chebroots),
        \texttt{cheby\_numba\_bisect.phi\_bisect} (1D bisection on the
        univariate polynomial), \texttt{cheby\_sym2.\{expand,contract\}}
        (symmetric subspace projection);
  \item Jacobian by forward difference at $\varepsilon = 10^{-6}$, dense
        $40\!\times\!40$, \texttt{numpy.linalg.eigvals}.
\end{itemize}

\end{document}
